% ============================================================
% Chapter 4: Debugging and Monitoring
% ============================================================
\section{调试与监控}

\subsection{调试 RabbitMQ 节点}

\subsubsection{初步启动}

\begin{figure}[H]
\centering
\includegraphics[width=0.35\textwidth]{image23}
\caption{RabbitMQ 初步启动}
\end{figure}

\subsubsection{停止}

\begin{lstlisting}
(rabbit@centos7-mq1)2> rabbit:stop().
ok
\end{lstlisting}

\begin{figure}[H]
\centering
\includegraphics[width=0.85\textwidth]{image24}
\caption{RabbitMQ 停止}
\end{figure}

\subsubsection{再次开启}

\begin{lstlisting}
(rabbit@centos7-mq1)4> rabbit:start().
ok
\end{lstlisting}

\subsection{节点位置无关性}

\begin{figure}[H]
\centering
\includegraphics[width=0.85\textwidth]{image25}
\caption{节点位置无关性 (1)}
\end{figure}

\begin{figure}[H]
\centering
\includegraphics[width=0.85\textwidth]{image26}
\caption{节点位置无关性 (2)}
\end{figure}

\subsection{调试加载系统库和当前目录}

\subsection{列出已加载的库（符号表）}

\begin{lstlisting}
code:all_loaded().
rp(code:all_loaded()).
\end{lstlisting}

\begin{figure}[H]
\centering
\includegraphics[width=0.85\textwidth]{image27}
\caption{已加载模块列表}
\end{figure}

\begin{itemize}[leftmargin=1.5em]
    \item 加载模块：\lstinline|l(ping).|
    \item 显示模块：\lstinline|{ping, "/home/ericksun/.../ping.beam"}|
\end{itemize}

\subsection{查看所有的路径}

使用 \lstinline|filename| 模块。

\subsection{唤起 xterm}

\begin{lstlisting}[style=shellstyle]
sudo yum install xterm.x86_64
\end{lstlisting}

\begin{figure}[H]
\centering
\includegraphics[width=0.85\textwidth]{image29}
\caption{xterm 界面}
\end{figure}

\subsection{查看 Monitor/Link 进程}

\begin{itemize}[leftmargin=1.5em]
    \item 查看 monitor 或 link 进程的情况
    \item 获取没有被 link 和 monitor 的进程
    \item 查看占用内存最大的进程
\end{itemize}

\subsection{日志}

\begin{itemize}[leftmargin=1.5em]
    \item \href{https://www.erlang.org/doc/man/logger.html#macros}{Erlang -- logger}
    \item \href{https://www.erlang.org/doc/man/erlang#type-reference}{Erlang -- erlang reference}
\end{itemize}

\subsection{再查清楚：日志 handler 是怎么注册的}

\subsubsection{旧版 error\_logger}

在较旧的 Erlang 版本里，\lstinline|error_logger| 是主要的日志和事件管理模块。可以借助 \lstinline|gen_event:which_handlers/1| 函数来查看 \lstinline|error_logger| 事件管理器里的所有处理程序。

\subsubsection{新版 logger 模块}

从 Erlang/OTP 21.0 版本开始，引入了新的 \lstinline|logger| 模块来替代 \lstinline|error_logger|。可以通过 \lstinline|logger:get_handler_config/0| 函数查看所有已配置的处理程序。

\subsection{查看 event 中所有的 handler}

\begin{lstlisting}
% Old version
gen_event:which_handlers(error_logger).

% New version (OTP 21+)
logger:get_handler_config().
\end{lstlisting}
